% TOP_PROBLEM: {"id": 3100089, "title": "Let G be a bicolored graph, and let H be the graph whose vertices are the valid pressing sequences of G and whose edges", "category_id": 2, "category": {"id": 2, "name": "combinatorics", "display_name": "Combinatorics", "description": "Counting problems, graph theory, discrete structures.", "slug": "combinatorics", "order_index": 2, "created_at": "2026-07-31T15:26:25.671Z"}, "status": "open"}
% FIRST_POSTED: 2026-10-11
% ENTRY_KIND: research_attempt
% SUBMISSION: {"kind":"research_attempt","category":"top_problems","problem_id":"3100089","model":{"provider":"OpenAI","version":"GPT 6.1 Sol","reasoning":"Max","generated_on":"2026-10-11"},"other_models":"Prior GPT-6-based Codex, exact version and reasoning unexposed, generated the preserved star draft. The same configured GPT 6.1 Sol Max generated the preserved double-star draft. These are process context, not dependencies of the counterexample proof. No subagents.","tools":"Public web browsing; Python 3.13.2 standard-library exhaustive search and direct edge/color and principal-determinant checks; NetworkX graph-atlas data (not the NetworkX package); local Git; repository structural validator; app LaTeX compiler. No theorem prover.","target":"For every finite simple black/white graph, any two successful pressing words can be connected through successful words at total insertion-deletion distance at most four; empty families satisfy the pairwise condition vacuously.","scope":"Full disproof of the original four-total-edit universal assertion: the explicit six-vertex graph in Theorem 1 has exactly two successful words at distance eight, hence H4 has two isolated vertices. This does not disprove the distinct four-deletions-from-each-word formulation.","human_contributions":"Under the guidance and supervision of Wentao Li: research direction, iterative prompting, proof/literature recheck requests and submission authorization. No recorded human mathematical proof development, substantial correction or independent verification.","verification":"Complete edge/color branching proof and block-order edit-distance proof rechecked on 2026-10-11; original verifier re-run; same-model adversarial audit. All 19 legal prefixes, 1957 distinct-symbol words, 64 principal subsets, LCS two and distance eight checked. No independent human or formal verification; historical novelty not established.","reproducibility":"Prompt A preceded discovery; complete Prompt B inputs, audits, corrected outputs and earlier drafts are retained locally outside the PR, without a public process archive. The complete manuscript and standalone repository verifier reproduce the counterexample. Presentation revisions preserve the proof and are recorded in Git history."}
% ATTEMPT_STATUS: solved
% FULL_SOLUTION_SCOPE: {"target":"The entire universal four-total-insertion/deletion connectivity statement of catalogue ID 3100089, with the original pressing rule and pairwise connectivity convention.","result":"Theorem 1: vertices 1..5 form a cycle, vertex 6 is adjacent only to 2, and initially exactly 4 and 5 are black. Its complete successful-word family is {432651,512643}, with insertion-deletion distance eight; H4 is disconnected.","coverage":"One admissible finite simple colored graph with a nonempty successful-word family disproves the universal assertion. The proof exhausts every legal branch for this graph and excludes every possible intermediate successful word. Other graph sizes, colorings, empty families and the empty graph need not be classified to disprove a universal statement.","dependencies":"No external theorem or computational certificate is needed for the self-contained proof; only the original pressing and insertion/deletion definitions. Literature and graph-atlas data supplied context and discovery.","human_verification":"None. Wentao Li provided direction and prompts; no independent human proof checking occurred.","unverified":"The complete proof has not been independently reviewed by a human or formalized in a theorem prover. Repository checks and same-model audits do not supply that verification.","novelty":"Expanded primary-source and duplicate search on 2026-10-11 located no matching colored five-cycle-with-pendant example in inspected material. Bixby-Flint-Miklos explicitly conjecture eight total edits, and this result does not refute that formulation. Whitlatch dissertation full text and JMM full proof remain inaccessible; historical novelty is unestablished."}
\documentclass[11pt]{article}
\usepackage[T1]{fontenc}
\usepackage{lmodern}
\usepackage[margin=25mm]{geometry}
\usepackage{amsmath,amssymb,amsthm,booktabs,microtype,xurl,hyperref}
\urlstyle{same}
\setlength{\emergencystretch}{2em}
\newtheorem{theorem}{Theorem}
\newtheorem{corollary}[theorem]{Corollary}
\newcommand{\ed}{\operatorname{ed}}
\newcommand{\LCS}{\operatorname{LCS}}
\title{A six-vertex counterexample to\\
four-total-edit pressing connectivity}
\author{OpenAI GPT 6.1 Sol\\\small Max reasoning}
\date{11 October 2026}
\AtBeginDocument{\maketitle}
\begin{document}
\begin{center}
\textbf{Abstract}
\end{center}
\begin{quote}
\small
We give a six-vertex bicolored graph with exactly two successful pressing
sequences. Their total insertion--deletion distance is eight, so the
four-total-edit connectivity assertion recorded as UnsolvedMath
ID 3100089 is false. The proof exhausts every legal branch directly.
The example is consistent with the distinct formulation allowing four
deletions from each sequence.
\end{quote}

\section{Introduction}
Successful pressing sequences of a bicolored graph need not form a
connected family under four total insertions and deletions. We prove this
by an explicit five-cycle with one pendant vertex. The result concerns
the literal four-total-edit statement in the repository definition of
UnsolvedMath ID 3100089 (AMR-030-0089) \cite{catalogue,cooper}; its
relationship to the eight-total-edit formulation is discussed in Section~3.

\subsection{Statement and conventions}
A bicolored graph is a finite simple graph whose vertices are black or
white. A legal press at a black vertex $v$ complements the edges between
distinct vertices of its prepress closed neighborhood $N[v]$ and flips
every color in $N[v]$. Other edges and colors are unchanged. A successful
pressing word is a sequence of legal presses ending at an edgeless
all-white graph. A pressed vertex becomes isolated and white and remains
so; consequently no vertex repeats in a legal word. Words retain the
individual vertex labels.

Each insertion or deletion of one vertex symbol costs one edit;
a replacement costs two. Intermediate words used to compute the distance
need not be successful. Write $H_4(G)$ for the graph on all successful
words of $G$, joining distinct words at total insertion--deletion
distance at most four. The recorded question is:
\begin{quote}
For every bicolored graph $G$, can every pair of successful pressing
words be joined by a finite path in $H_4(G)$?
\end{quote}
The pairwise condition is vacuous for an empty family and holds for a
singleton. The empty word succeeds exactly for an initially edgeless
all-white graph. These conventions agree with the recorded target.

\section{Main result}
Let $G$ have vertices $V=\{1,2,3,4,5,6\}$, edges
\[
 E=\bigl\{\{1,2\},\{2,3\},\{3,4\},\{4,5\},\{5,1\},\{2,6\}\bigr\},
 \qquad B=\{4,5\},
\]
where $B$ is the set of initially black vertices. Thus $1,2,3,4,5$ form a
five-cycle and $6$ is a white pendant vertex at $2$. This is a connected
finite simple graph and has black vertices, so the counterexample does
not rely on vacuous connectivity or on lack of successful pressing words.

\begin{center}
\begin{tabular}{lll}
\toprule
Vertex & Neighbors & Initial color\\
\midrule
$1$ & $2,5$ & white\\
$2$ & $1,3,6$ & white\\
$3$ & $2,4$ & white\\
$4$ & $3,5$ & black\\
$5$ & $1,4$ & black\\
$6$ & $2$ & white\\
\bottomrule
\end{tabular}
\end{center}

\begin{theorem}\label{counterexample}
The complete successful-word family of $G$ is
\[
 \sigma=4\,3\,2\,6\,5\,1,
 \qquad
 \tau=5\,1\,2\,6\,4\,3.
\]
Their total insertion--deletion distance is eight. Consequently $H_4(G)$
consists of two isolated vertices, and the original universal connectivity
assertion is false.
\end{theorem}
\begin{proof}
We first determine every successful word directly from the graph rule.
Write $ij$ for the unordered edge $\{i,j\}$, and write a pressing prefix
without separators since every symbol is a single digit. The color-preserving
automorphism
\[
             \rho=(1\ 3)(4\ 5),\qquad \rho(2)=2,\quad\rho(6)=6
\]
exchanges the two initially black vertices. Every legal nonempty word
starts with $4$ or $5$. It therefore suffices to analyze the branch
starting with $4$, and then apply $\rho$.

The complete half-tree is the following table. All six vertices remain
present; vertices already pressed are isolated and white. The right
column lists \emph{all} available legal next presses.
\begin{center}
\begin{tabular}{lll}
\toprule
Prefix & Remaining edges & Black vertices\\
\midrule
$\varnothing$ & $12,15,23,26,34,45$ & $4,5$\\
$4$ & $12,15,23,26,35$ & $3$\\
$43$ & $12,15,25,26$ & $2,5$\\
$435$ & $26$ & $1$\\
$4351$ & $26$ & none\\
$432$ & $16,56$ & $1,6$\\
$4321$ & $56$ & none\\
$4326$ & $15$ & $5$\\
$43265$ & none & $1$\\
$432651$ & none & none\\
\bottomrule
\end{tabular}
\end{center}

Here is also an explicit derivation of the branching rows. Pressing $4$
uses $N[4]=\{3,4,5\}$: the edges $34,45$ disappear, $35$ appears,
and only $3$ is black. Pressing that forced $3$ uses
$\{2,3,5\}$: $23,35$ disappear, $25$ appears, and the black set becomes
$\{2,5\}$. The remaining edges are $12,15,25,26$.

If $5$ is chosen next, its closed neighborhood is $\{1,2,5\}$.
All three edges of that triangle disappear, leaving only $26$ and the
black vertex $1$. Pressing the isolated $1$ leaves an all-white graph
with edge $26$, so this branch cannot succeed.

If $2$ is chosen after $43$, its closed neighborhood is
$\{1,2,5,6\}$. All edges incident with $2$ disappear; $15$ disappears
and $16,56$ appear. The black set becomes $\{1,6\}$. Pressing $1$
at this point removes $16$ and makes $6$ white, leaving the all-white
edge $56$; this branch also fails. Choosing $6$ instead uses
$\{1,5,6\}$, removes $16,56$, creates $15$, and leaves $5$ as the
only black vertex. The next press $5$ removes $15$ and makes $1$
black; pressing $1$ then gives the edgeless all-white graph.

All decisions have been exhausted: the only successful word starting
with $4$ is $\sigma$. Applying $\rho$ gives the only successful word
starting with $5$, namely $\tau$. The initial graph is not terminal,
so the empty word adds no success. A stalled all-white graph admits
no further legal press, and a successful terminal state also admits
no further press. This proves completeness of the two-word family,
without any matrix theorem or computational assumption.

For the distance, divide the symbols into the ordered pairs
\[
                 P=(4,3),\qquad Q=(2,6),\qquad R=(5,1).
\]
The words have block orders $\sigma=PQR$ and $\tau=RQP$, with the
same internal order within each pair. The relative order of any two
symbols from different blocks is reversed between the two words.
Thus a common subsequence can contain symbols from only one block,
and has length at most two. Each pair itself is common, so
$\LCS(\sigma,\tau)=2$.

Under insertions and deletions, the original symbols that survive
form a common subsequence. At most two of the six can survive, so
at least four deletions and four insertions are necessary. Conversely,
keep the common pair $26$, delete the other four symbols from
$\sigma$, and insert $51$ before it and $43$ after it. This gives
$\tau$ with eight edits. Hence $\ed(\sigma,\tau)=8$ exactly.
Since these are the only successful words, they have no edge in
$H_4(G)$ and there is no intermediate successful vertex through which
to connect them. This disproves the universal statement.
\end{proof}

\begin{corollary}\label{lowerbound}
Any total insertion--deletion threshold that connects successful words
for every bicolored graph must be at least eight. For the displayed graph,
eight is exactly the smallest threshold giving connectivity.
\end{corollary}
\begin{proof}
The same two-vertex word graph has no edge at any threshold below eight,
and has its single connecting edge at threshold eight.
\end{proof}

\section{Relation to the pressing game conjecture}
Cooper--Davis \cite[Conjecture 1]{cooper} and the repository definition
use the four-insertions-or-deletions wording considered here. In contrast,
Bixby--Flint--Mikl\'os \cite[Conjecture 6]{bixby} permit a
longest-common-subsequence deficit of four. For two words of common
length $k$ and LCS length $k-d$, the total insertion--deletion distance
is $2d$; their formulation thus permits eight total edits. The two words
in Theorem~\ref{counterexample} have deficit four and meet that bound.
The stronger path result in \cite[Theorem 7]{bixby} is also consistent
with this example, whose graph is not a path.

The indexed JMM 2024 abstract \cite{jmm} announces connectivity with
at most four deletions from each neighboring word, again the
eight-total-edit formulation. Its underlying proof was not accessed.
Our result disproves the literal four-total-edit assertion; it does not
refute Bixby--Flint--Mikl\'os Conjecture 6 or the JMM announcement.
The literature-search limits are recorded in Appendix~B.

\section{Completion Estimate}
\noindent\textbf{COMPLETION: 100\%}
Theorem~\ref{counterexample} disproves the entire recorded universal
four-total-edit assertion. This is claimed coverage of the original
target, not a measure of confidence, novelty, or independent verification.

\section*{Appendix A. Verification and contributions}
\subsection*{A.1 Computational verification}
The standalone verifier enumerates the entire legal tree: 19 prefixes,
the two successful words, and the four maximal stalled words
$4321,4351,5123,5143$. It checks every proof-table row and the reflection.
Dense principal determinants over $\mathbb F_2$ check all 64 vertex
subsets and all 1,957 distinct-symbol words, including 720 full
permutations; the legal flags agree with the direct tree.
Edit dynamic programming and common-subsequence enumeration confirm
distance eight and LCS length two. A separate literal-rule replay
agrees on every legal state. These checks passed with Python 3.13.2.

Run from the repository root with assertions enabled:
\begin{verbatim}
python3 scripts/verify_pressing_3100089_counterexample.py \
  --output /absolute/path/to/counterexample-results.json
\end{verbatim}
The JSON records the graph, complete state tree, outcomes, counts,
distance, software version, and verifier-source SHA-256.
The proof itself requires no computational certificate.

\subsection*{A.2 Contributions and scope}
\paragraph{Contributor declaration.}
Under the guidance and supervision of Wentao Li, OpenAI GPT 6.1 Sol,
Max reasoning generated the search, counterexample, proof, verifier,
manuscript, and model-assisted audits on 11 October 2026. Wentao Li
provided research direction and iterative prompts, directed the proof
and literature rechecks, and authorized submission. No human mathematical
proof development, substantial mathematical correction, or independent
human proof verification is recorded.

\paragraph{Full-solution scope declaration.}
Theorem~\ref{counterexample} refutes the entire universal four-total-edit
target: one admissible colored graph has a complete successful-word
family with disconnected $H_4$. Every legal branch is exhausted;
other graph sizes, colorings, and empty families need not be classified
to disprove the universal quantifier. There are no external theorem
dependencies. The complete proof has not been independently reviewed
by a human or formalized in a theorem prover. The same model performed
the audits and generated the checking programs. Historical novelty is
unestablished. The manuscript claims no resolution of the distinct
universal eight-total-edit formulation. The repository records
\texttt{ATTEMPT\_STATUS: solved} as an LLM claim; source status remains
\emph{open}, and community review remains separate.

\subsection*{A.3 Process and reproducibility}
Discovery used standard-library Python enumeration and public NetworkX
graph-atlas data \cite{atlas}, without importing NetworkX. The search
found atlas index 104 with black-color mask 24 after 4,596 graph/color
cases and 97,571 successful words. It was not an exhaustive atlas search,
and no vertex-minimality claim is made.

Prompt A preceded discovery. Complete Prompt B inputs, audits, corrected
outputs, source checks, and earlier drafts are retained locally in the
dated \texttt{research/} folders outside the contribution repository,
without a public process archive. Git history preserves the submitted
version and its presentation revisions; the proof is unchanged.
The earlier star draft used GPT-6-based Codex with exact version and
reasoning unexposed; the double-star draft used the same configured
GPT 6.1 Sol Max. These drafts are process context, not proof dependencies.
No other model, subagent, or theorem prover was used. Public web browsing,
Python, Git, repository validation, and the app LaTeX compiler supported
the work; formatting and build checks do not verify mathematics.

\section*{Appendix B. Literature-search record}
The search on 11 October 2026 inspected the Cooper--Davis and
Bixby--Flint--Mikl\'os papers, the accessible uniquely pressable graph
and poset papers \cite{unique,vposets}, the author poster \cite{poster},
Mikl\'os's dissertation \cite[Chapter 9, pp.~167--178]{miklosthesis},
Cooper's problems page \cite{problems}, and the indexed JMM abstract.
The poset paper's four-vertex examples and the poster's six-vertex graph
are distinct from the present construction. Searches covered pressing
counterexamples, two successful words, six-vertex graphs, pendant cycles,
and equivalent edit/LCS formulations. Pre-submission repository issue
and PR searches for the ID and pressing terminology found no match.

Whitlatch's dissertation landing page and abstract \cite{whitlatchthesis}
were accessible, but its linked full text returned HTTP 403. The JMM
underlying proof was not accessed. No matching example was located in
the inspected material; this does not establish historical novelty or
exclude an equivalent earlier construction. None of these sources is
a dependency of the self-contained proof.

\begin{thebibliography}{99}
\bibitem{cooper} J. Cooper and J. Davis,
\emph{Successful Pressing Sequences for a Bicolored Graph and Binary Matrices},
arXiv:1502.07450v5, 2015, Definitions 1--2 and Conjecture 1, pp.~2--4.
\url{https://arxiv.org/abs/1502.07450v5}.
\bibitem{bixby} E. Bixby, T. Flint, and I. Mikl\'os,
\emph{Proving the pressing game conjecture on linear graphs},
arXiv:1303.6799, 2013, Conjecture 6 and Theorem 7;
Involve 9 (2016), 41--56.
\url{https://arxiv.org/abs/1303.6799}.
\url{https://doi.org/10.2140/involve.2016.9.41}.
\bibitem{jmm} B. Li, I. Mikl\'os, and C. R. Teplica,
\emph{The Connectedness of the Black-and-White Graph Pressing Game},
Joint Mathematics Meetings, San Francisco, 3--6 January 2024,
abstract 1192-05-32487, received 12 September 2023.
Indexed abstract consulted; full proof not accessed.
\url{https://www.ams.org/journals/abs/2024-45-01/1192.pdf}.
\bibitem{problems} J. Cooper, \emph{Combinatorial Problems I Like},
pressing-sequence connectivity entry, page last edited October 2020;
consulted 11 October 2026.
\url{https://people.math.sc.edu/cooper/combprob.html}.
\bibitem{unique} J. N. Cooper and H. W. Whitlatch,
\emph{Uniquely Pressable Graphs: Characterization, Enumeration, and Recognition},
arXiv:1706.07468. Accessible text consulted 11 October 2026.
\url{https://arxiv.org/abs/1706.07468}.
\bibitem{vposets} J. Cooper, P. Gartland, and H. Whitlatch,
\emph{A New Characterization of $\mathcal V$-Posets},
arXiv:1810.07276. Accessible text and Figure 2 inspected 11 October 2026.
\url{https://arxiv.org/abs/1810.07276}.
\bibitem{poster} J. Cooper and H. Whitlatch,
\emph{Genome Inversions and Graph Pressing Sequences}, research poster,
public author-hosted PDF, consulted 11 October 2026.
\url{https://connect.gonzaga.edu/asset/file/1363/Clemson_Poster.pdf}.
\bibitem{miklosthesis} I. Mikl\'os,
\emph{Computational complexity of counting and sampling problems in bioinformatics},
dissertation dc\_2046\_22, R\'enyi Institute, Chapter 9.
Conjecture statement and discussion inspected 11 October 2026.
\url{https://real-d.mtak.hu/1570/7/dc_2046_22_doktori_mu.pdf}.
\bibitem{whitlatchthesis} H. W. Whitlatch,
\emph{Successful Pressing Sequences in Simple Pseudo-Graphs},
Ph.D. dissertation, University of South Carolina, Spring 2019.
Landing page and abstract consulted; linked full text inaccessible.
\url{https://scholarcommons.sc.edu/etd/5135/}.
\bibitem{atlas} NetworkX developers, graph atlas data and documentation,
consulted 11 October 2026; based on R. C. Read and R. J. Wilson,
\emph{An Atlas of Graphs}, Oxford University Press, 1998.
Discovery data only; not a dependency of the disproof.
\url{https://networkx.org/documentation/stable/reference/generated/networkx.generators.atlas.graph_atlas_g.html}.
\url{https://github.com/networkx/networkx/blob/main/networkx/generators/atlas.dat.gz}.
\bibitem{catalogue} ULAM.AI and the UnsolvedMath Contributors,
\emph{UnsolvedMath}, record 3100089 (AMR-030-0089), and the repository
definition at commit b1f633121074e16579fa7fd05314e5f504d29bae.
Original curation and metadata: CC BY 4.0.
\url{https://huggingface.co/datasets/ulamai/UnsolvedMath}.
\url{https://github.com/mehmetmars7/Erdosproblems-llm-hunter/blob/b1f633121074e16579fa7fd05314e5f504d29bae/attacks/open_problems/top_problems/definitions/3100089.tex}.
\url{https://creativecommons.org/licenses/by/4.0/}.
\end{thebibliography}
\end{document}
